\changecolours{ScoutGreen}
\section[Anatomia da Relação]{Domínios, Atributos, Tuplas, Grau e Cardinalidade}

% ------------------------------------------------------------------------------
\begin{frame}{Dicionário de Equivalência Terminológica e Anatomia da Relação}
  \vspace{-0.1cm}
  \begin{columns}[t]
    \begin{column}{0.46\textwidth}
      \alert{Três Vocabulários Equivalentes}\par
      \vspace{0.1cm}
      \scriptsize
      \begin{tabular}{lll}
        \toprule
        \textbf{Formal (Codd)} & \textbf{Tabular / Lógico} & \textbf{Arquivos} \\
        \midrule
        Relação ($R$) & Tabela & Arquivo \\
        Tupla ($t$) & Linha & Registro \\
        Atributo ($A_i$) & Coluna & Campo \\
        Domínio ($D_i$) & Tipo / Formato & Máscara \\
        Grau ($n$) & Nº de Colunas & Nº Campos \\
        Cardinalidade ($m$) & Nº de Linhas & Nº Registros \\
        \bottomrule
      \end{tabular}
      \vspace{0.15cm}\par
      $\bullet$ \textbf{Grau (Aridade):} Propriedade \textbf{estática} do esquema $R$.\\
      $\bullet$ \textbf{Cardinalidade:} Propriedade \textbf{dinâmica} do estado $r(R)$ (varia conforme novas tuplas entram ou saem).
    \end{column}
    \begin{column}{0.51\textwidth}
      \alert{Visualização Formal: Relação \texttt{MEDICO}}\par
      \vspace{0.1cm}
      \begin{center}
        \tikz[scale=0.68, transform shape]{
          \node[font=\bfseries\sffamily\scriptsize, ScoutNavy] at (2.7, 2.1) {Esquema $R$: \texttt{MEDICO}(\underline{matricula}, nome, crm, cod\_depto)};

          % Cabecalho (Atributos)
          \draw[fill=ScoutGreen!20, draw=ScoutGreen, thick] (0, 1.1) rectangle (5.4, 1.7);
          \node[font=\bfseries\sffamily\tiny] at (0.65, 1.4) {\underline{matricula}};
          \node[font=\bfseries\sffamily\tiny] at (2.05, 1.4) {nome};
          \node[font=\bfseries\sffamily\tiny] at (3.45, 1.4) {crm};
          \node[font=\bfseries\sffamily\tiny] at (4.75, 1.4) {cod\_depto};

          % Tuplas
          \draw[fill=white, draw=black!50] (0, 0.5) rectangle (5.4, 1.1);
          \node[font=\sffamily\tiny] at (0.65, 0.8) {1001};
          \node[font=\sffamily\tiny] at (2.05, 0.8) {Dra. Ana};
          \node[font=\sffamily\tiny] at (3.45, 0.8) {CE-1234};
          \node[font=\sffamily\tiny] at (4.75, 0.8) {10};

          \draw[fill=black!3, draw=black!50] (0, -0.1) rectangle (5.4, 0.5);
          \node[font=\sffamily\tiny] at (0.65, 0.2) {1002};
          \node[font=\sffamily\tiny] at (2.05, 0.2) {Dr. Bruno};
          \node[font=\sffamily\tiny] at (3.45, 0.2) {CE-5678};
          \node[font=\sffamily\tiny] at (4.75, 0.2) {10};

          \draw[fill=white, draw=black!50] (0, -0.7) rectangle (5.4, -0.1);
          \node[font=\sffamily\tiny] at (0.65, -0.4) {1003};
          \node[font=\sffamily\tiny] at (2.05, -0.4) {Dra. Carla};
          \node[font=\sffamily\tiny] at (3.45, -0.4) {CE-9012};
          \node[font=\sffamily\tiny\itshape, gray] at (4.75, -0.4) {NULO};

          % Cotas Grau e Cardinalidade
          \draw[<->, thick, ScoutNavy] (0, -1.1) -- node[below, font=\tiny\bfseries] {Grau (Aridade) $n = 4$ atributos} (5.4, -1.1);
          \draw[<->, thick, ScoutGreen] (5.7, -0.7) -- node[right, font=\tiny\bfseries, align=left] {Cardinalidade\\$m = 3$ tuplas} (5.7, 1.1);
        }
      \end{center}
    \end{column}
  \end{columns}
\end{frame}

% ------------------------------------------------------------------------------
\twocolframe[title={Definição Matemática de Esquema ($R$) e Instância ($r(R)$)},leftcol=7.0cm,rightcol=7.0cm]{%
  \alert{Domínio Atômico e Esquema Intensional}\par
  \itemseps[topsepi=0.5mm,itemsepi=1mm,labelsep=2mm,leftmargini=3mm]
  \begin{itemize}
    \item \textbf{Domínio ($D_i$):} Conjunto de valores \textbf{atômicos} (indivisíveis) permitidos para um atributo, associado a uma definição semântica e formato.
    \item \textbf{Esquema de Relação ($R$):} Denotado por $R(A_1, A_2, \dots, A_n)$, define a estrutura intensional (cabeçalho invariante) composta por $n$ atributos onde $\text{dom}(A_i) = D_i$.
  \end{itemize}
}{%
  \alert{Estado/Instância como Subconjunto Cartesiano}\par
  \itemseps[topsepi=0.5mm,itemsepi=1mm,labelsep=2mm,leftmargini=3mm]
  \begin{itemize}
    \item \textbf{Instância Extensional $r(R)$:} É um conjunto finito de $m$ tuplas $\{t_1, t_2, \dots, t_m\}$ que reflete o estado dos dados em dado instante:
      $$r(R) \subseteq \text{dom}(A_1) \times \text{dom}(A_2) \times \dots \times \text{dom}(A_n)$$
    \item \textbf{Cada Tupla $t$:} É uma lista ordenada de $n$ valores $t = \langle v_1, v_2, \dots, v_n \rangle$, onde cada $v_i \in \text{dom}(A_i)$ ou é um valor nulo especial (\textit{NULL}).
  \end{itemize}
}
